\section{Introduction}
Autoregressive (AR) Transformers~\citep{vaswani2017attention,radford2019language,brown2020language} generate one token per network evaluation, so latency grows linearly with output length. Two families of parallel generators have been proposed to remove this bottleneck. Non-autoregressive and iterative-refinement models~\citep{gu2018nat,lee2018deterministic,ghazvininejad2019maskpredict,gu2019levenshtein} and, more recently, discrete diffusion language models (DLMs)~\citep{austin2021d3pm,li2022diffusionlm,lou2024sedd,sahoo2024mdlm,shi2024md4,gulrajani2023plaid,nie2025llada} predict many tokens per step, but they revisit the full sequence at every step and degrade quickly when the number of steps is reduced. Next-Concept Modeling (NCM)~\citep{kumbong2026ncm} takes a third route that is closer to next-scale prediction in vision~\citep{tian2024var}: a multiscale residual VQ-VAE (MS-RVQ) encodes a sequence into $K$ levels of discrete codes of increasing length, and a Transformer generates one level per forward pass, emitting all codes within a level in parallel. On the released checkpoints, NCM solves 99.4\% of easy Sudoku boards with 12 network evaluations and is 10--31$\times$ faster than diffusion baselines on text~\citep{kumbong2026ncm}.

Two structural properties of NCM motivate this work. First, the latent state is \emph{purely additive}: the cumulative reconstruction after level $k$ is the sum of the upsampled code embeddings of levels $0,\dots,k$, and nothing in the sampler can revise a coarse code once it has been emitted. A wrong code at a coarse level, where a single position covers many output tokens, is broadcast over the whole sequence by linear interpolation, and finer codebooks, which were fitted to \emph{clean} residuals, are evaluated out of distribution. Second, the depth of generation is \emph{fixed}: every input pays for all $K$ levels, although the official tokenizers are trained with fine-level dropout (\texttt{fine\_drop}), a nested-dropout scheme~\citep{rippel2014nested} that makes intermediate cumulative latents decodable by the frozen decoder.

These observations suggest two training-free modifications borrowed from continuous-time samplers and adaptive-computation networks: a learned \emph{corrector} that revises the cumulative latent after every predictor step, as in predictor-corrector diffusion samplers~\citep{song2021sde,zhao2023unipc}, and a \emph{halting rule} that stops generation once the latent has converged, as in adaptive computation time and early-exit Transformers~\citep{graves2016act,banino2021pondernet,elbayad2020depth,schuster2022calm}. We call the combination Adaptive Predictor-Corrector NCM (APC-NCM). The central question of this paper is: \emph{can predictor-corrector latent revision and adaptive early halting improve a frozen coarse-to-fine NCM without retraining it?} We answer it in two parts. We first derive the conditions each mechanism requires, then test these conditions on the released checkpoints under a strict protocol: official code, no access to targets at inference, hyper-parameters selected only on held-out inputs, five seeds, and paired significance tests. The answer is negative, and the analysis identifies why.

Our contributions are:
\begin{itemize}
\item \textbf{Formulation.} We write the NCM sampler as a predictor that adds one quantized residual per level, formalize the additive-drift failure mode, and define a learned latent corrector $C_\omega$ applied in the loop, on the decode path, or with projection back onto the codebook lattice, as well as batch-level and per-sequence halting rules.
\item \textbf{Conditions for success.} Proposition~\ref{prop:contraction} shows that additive MS-RVQ error grows linearly in the per-level code error, whereas a corrector that is $\epsilon_{\mathrm{corr}}$-accurate and contractive along the \emph{closed-loop} trajectory yields a bounded error governed by $\prod_j(1-\gamma_j)$. Proposition~\ref{prop:rd} bounds the KL cost of halting by $L_D^2\tau^2/(2(1-\rho)^2)$ under geometric increment decay $\rho<1$, and a Pinsker argument turns this into exact output preservation when the final decoding has sufficient margin.
\item \textbf{Bit-exact controlled evaluation.} We reproduce the official baselines on a T4 (Sudoku easy/medium/hard: \OffEasy/\OffMedium/\OffHard\% solved; WMT14 De--En: \OffWmt{} BLEU; XSum: \OffXsumRL{} ROUGE-L) with a harness that wraps the official sampler unchanged, and evaluate every variant with 5 seeds, Wilcoxon and McNemar tests, and measured FLOPs, latency and memory.
\item \textbf{Three diagnosed violations.} (i) \emph{Off-manifold exposure bias}: the generator is brittle to continuous latents that are not sums of codewords; correction lowers open-loop error but collapses closed-loop generation (Sudoku-hard \LoopVoneHard\% vs.\ \BaseHard\%), validation selection drives the step size to zero, codebook projection recovers most of the loss without beating NCM, and DAgger fine-tuning of the generator yields no significant gain. (ii) \emph{Non-decaying residual flux}: fine-drop makes intermediate latents decodable but does not make increments contract ($\hat\rho>1$ at most transitions), so Proposition~\ref{prop:rd} is vacuous. (iii) \emph{Terminal-level concentration}: task resolution occurs at level $K$, so no lossless early exit exists even per sequence.
\item \textbf{Prescriptions.} We derive two pre-training changes that would satisfy the propositions: a residual-decay penalty in Stage~1 and continuous latent perturbations in Stage~2.
\end{itemize}

\section{Related Work}
\textbf{Coarse-to-fine and next-scale generation.} VAR~\citep{tian2024var} generates images as a sequence of token maps of increasing resolution, each predicted in parallel; NCM~\citep{kumbong2026ncm} transfers the idea to 1-D text with a block-diagonal attention pattern and a frozen multiscale tokenizer. Earlier hierarchical latent models such as VQ-VAE-2~\citep{razavi2019vqvae2} and masked generative transformers such as MaskGIT~\citep{chang2022maskgit} also decouple global structure from detail. None of these revise an already emitted coarse code.

\textbf{Diffusion and iterative-refinement language models.} Discrete diffusion~\citep{austin2021d3pm,campbell2022ctmc,lou2024sedd,sahoo2024mdlm,shi2024md4}, continuous embedding diffusion~\citep{li2022diffusionlm,gulrajani2023plaid} and large masked diffusion LMs~\citep{nie2025llada} refine the full sequence at every step. Non-autoregressive translation~\citep{gu2018nat}, iterative refinement~\citep{lee2018deterministic}, Mask-Predict~\citep{ghazvininejad2019maskpredict} and the Levenshtein Transformer~\citep{gu2019levenshtein} edit previous hypotheses, which is the discrete analogue of a corrector.

\textbf{Residual vector quantization.} Multi-stage VQ~\citep{juang1982multiple} and its neural descendants~\citep{vandenoord2017vqvae,lee2022rqvae,zeghidour2022soundstream,defossez2023encodec} represent a vector as a sum of codewords chosen greedily on the running residual. Structured dropout of trailing stages~\citep{rippel2014nested,zeghidour2022soundstream} makes truncated sums decodable, which is the property the NCM \texttt{fine\_drop} auxiliary exploits.

\textbf{Adaptive computation and early exit.} ACT~\citep{graves2016act}, Universal Transformers~\citep{dehghani2019universal} and PonderNet~\citep{banino2021pondernet} learn how many recurrent steps to spend; BranchyNet~\citep{teerapittayanon2016branchynet}, DeeBERT~\citep{xin2020deebert}, depth-adaptive Transformers~\citep{elbayad2020depth} and CALM~\citep{schuster2022calm} exit from intermediate layers when a confidence or state-change criterion is met. Our halting rules are of the state-change type, applied to generation levels rather than layers.

\textbf{Predictor-corrector samplers and exposure bias.} Predictor-corrector schemes for score-based diffusion~\citep{ho2020ddpm,song2021sde} interleave a numerical integration step with Langevin correction; DDIM~\citep{song2021ddim}, DPM-Solver~\citep{lu2022dpmsolver} and UniPC~\citep{zhao2023unipc} design higher-order predictor and corrector updates. In sequence models, the mismatch between training on clean prefixes and testing on model-generated ones is known as exposure bias, studied through scheduled sampling~\citep{bengio2015scheduled} and DAgger~\citep{ross2011dagger}. Our negative result for in-loop correction is an instance of this mismatch.

\textbf{Constraint reasoning benchmarks.} Sudoku has been used to probe relational and constraint reasoning in neural networks~\citep{palm2018rrn,wang2019satnet}. We follow NCM and do not inject any rule-based constraint into the sampler.

\section{Preliminaries: Next-Concept Modeling}
\label{sec:prelim}
\subsection{Stage 1: multiscale residual VQ-VAE}
Let $x\in\mathcal{X}^{T}$ be a pre-tokenized sequence and $z=E_\theta(x)\in\mathbb{R}^{L\times d}$ the continuous encoder output ($L=T$ in all released models since the compression factor is 1). The MS-RVQ uses $K$ levels with lengths $\ell_0<\dots<\ell_{K-1}=L$ and codebooks $\mathcal{V}_k=\{e_{k,v}\}_{v=1}^{|\mathcal{V}_k|}\subset\mathbb{R}^d$. With $\mathrm{Down}_k$ the area interpolation to $\ell_k$ positions (a learned strided convolution at $k=0$) and $\Up_k$ the linear interpolation back to $L$ (a learned transposed convolution at $k=0$), the encoder computes for $k=0,\dots,K-1$
\begin{align}
r^\ast_{k} &= \argmax_{v}\; -\big\|\mathrm{Down}_k\big(z-\zh^{(k-1)}_{\mathrm{GT}}\big)-e_{k,v}\big\|_2^2, \label{eq:enc}\\
\zh^{(k)}_{\mathrm{GT}} &= \zh^{(k-1)}_{\mathrm{GT}} + \Phi_k\!\big(\Up_k(E_k(r^\ast_k))\big), \label{eq:cum}
\end{align}
with $\zh^{(-1)}_{\mathrm{GT}}=0$, where $E_k(\cdot)$ looks up codewords and $\Phi_k(h)=(1-\varrho)h+\varrho\,\mathrm{Conv1d}_k(h)$ with $\varrho=0.5$ is the residual refinement (\texttt{quant\_resi}). The argmin is taken row-wise and the arg-max form in~\eqref{eq:enc} is used only for notational uniformity. In the code, \eqref{eq:enc}--\eqref{eq:cum} are \texttt{MultiscaleResidualVectorQuantizer.z\_to\_idx}, which also returns the list $(\zh^{(0)}_{\mathrm{GT}},\dots,\zh^{(K-1)}_{\mathrm{GT}})$ when \texttt{return\_zhats=True}. The decoder $D_\psi$ maps $\zh^{(K-1)}$ to per-position logits.

\textbf{Fine-level dropout.} During tokenizer training, with probability $p_{\mathrm{fd}}$ ($0.5$ for Sudoku, WMT14, XSum and TinyStories) the reconstruction passed to $D_\psi$ is truncated to $\zh^{(J)}$ with $J\sim\mathcal{U}\{k_{\min}-1,\dots,K-2\}$. The decoder is therefore trained on the mixture
\begin{equation}
\pi_{\mathrm{fd}}(J) = (1-p_{\mathrm{fd}})\,\delta_{K-1}(J) + \tfrac{p_{\mathrm{fd}}}{K-k_{\min}}\textstyle\sum_{j=k_{\min}-1}^{K-2}\delta_j(J), \label{eq:finedrop}
\end{equation}
so intermediate cumulative latents lie on the decoder's training support.

\subsection{Stage 2: next-concept Transformer}
The generator $G_\phi$ is a Transformer with a prefix of $P$ conditioning tokens (puzzle, source sentence or article) and block-diagonal attention and rotary position embeddings~\citep{su2024roformer}: level $k+1$ attends to the prefix and to its own $\ell_{k+1}$ positions. At generation level $k$ the official loop computes
\begin{align}
\zh^{(k)} &= \zh^{(k-1)} + \Phi_k\big(\Up_k(E_k(\hat r_k))\big), \label{eq:pred}\\
u_{k+1} &= \mathrm{Down}_{k+1}\big(\zh^{(k)}\big)\in\mathbb{R}^{\ell_{k+1}\times d},\label{eq:down}\\
\hat r_{k+1} &= \argmax\; G_\phi\big(u_{k+1},\mathrm{prefix}\big)\in\{1,\dots,|\mathcal{V}_{k+1}|\}^{\ell_{k+1}}, \label{eq:gen}
\end{align}
which corresponds to \texttt{get\_next\_autoregressive\_input} in \texttt{models/quant.py} and the level loop of \texttt{NextConceptTransformer.generate} in \texttt{models/next\_concept.py}. All released configurations decode greedily (temperature 0 or top-$k=1$); XSum additionally applies the official frequency penalty when reading tokens from $D_\psi$. One forward pass of $G_\phi$ is spent per level, and the output is $D_\psi(\zh^{(K-1)})$.

\section{Proposed Methodology: APC-NCM}
\subsection{Additive drift}
Define the per-level code error $\epsilon_k=\Phi_k\Up_k\big(E_k(\hat r_k)-E_k(r^\ast_k)\big)$, where $r^\ast_k$ are the codes the tokenizer assigns to the reference output. Subtracting~\eqref{eq:cum} from~\eqref{eq:pred} gives the drift recursion
\begin{equation}
e^{(k)}\triangleq\zh^{(k)}-\zh^{(k)}_{\mathrm{GT}} = e^{(k-1)}+\epsilon_k=\sum_{j=0}^{k}\epsilon_j. \label{eq:drift}
\end{equation}
Because $\Up_j$ is a linear interpolation from $\ell_j\ll L$ positions, a single wrong code at a coarse level changes $\lceil L/\ell_j\rceil$ output positions. Moreover, $\hat r_{k+1}$ is computed from $u_{k+1}$, which contains $e^{(k)}$, while the generator and codebook $\mathcal{V}_{k+1}$ were trained on $u_{k+1}$ built from clean residuals, so nothing in the sampler is designed to contract~\eqref{eq:drift}; whether finer levels nevertheless cancel drift is an empirical question (Remark~\ref{rem:c0}).

\subsection{Learned latent corrector $C_\omega$}
We insert a learned revision after the predictor update~\eqref{eq:pred}:
\begin{equation}
\zt^{(k)}=\zh^{(k)}_{\mathrm{pred}}-\eta_k\,C_\omega\big(\zh^{(k)}_{\mathrm{pred}},h_{\mathrm{prefix}},k\big), \label{eq:corr}
\end{equation}
with $\eta_k=s_k\,\sigma(\beta_k)\in[0,1]$, where $\sigma$ is the logistic function, $\beta_k$ a learnable logit and $s_k\in\{0,0.25,0.5,1\}$ a scale chosen on validation data (Section~\ref{sec:select}). $h_{\mathrm{prefix}}\in\mathbb{R}^{P\times d}$ is the frozen raw-token prefix embedding that $G_\phi$ itself uses (\texttt{raw\_token\_prefix}), so $C_\omega$ sees exactly the information available to the generator. The network normalizes its input by its RMS $\varsigma=\|\zh\|_F/\sqrt{Ld}$, projects to width $h$, adds a level embedding, applies two depthwise-separable 1-D convolution blocks (kernel 5)~\citep{chollet2017xception}, one cross-attention block over $h_{\mathrm{prefix}}$ (4 heads) and an MLP, and maps back to $\mathbb{R}^{L\times d}$ through a zero-initialized projection whose output is rescaled by $\varsigma$. Zero initialization makes~\eqref{eq:corr} the identity at the start of training. Widths $h\in\{140,316,516\}$ give $0.31$M, $1.23$M and $3.04$M parameters on Sudoku ($d=256$).

\textbf{Two insertion points.} (i) \emph{In-loop} (the original proposal): $\zt^{(k)}$ replaces $\zh^{(k)}$ in~\eqref{eq:down}, so it changes all later generator inputs. (ii) \emph{Decode-path}: the generator continues from the uncorrected $\zh^{(k)}$, and $\zt^{(k)}$ is used only when a latent is decoded (at a halting level or at the final level). The decode-path variant cannot shift the generator's input distribution; it is trained for all $K$ levels.

\textbf{Training objective.} Generator and tokenizer are frozen. For a training input with reference $y$, the targets $\zh^{(k)}_{\mathrm{GT}}$ come from \texttt{z\_to\_idx}$(E_\theta(y))$. Predicted trajectories $\zh^{(k)}_{\mathrm{pred}}$ come from two sources: (a) free-running greedy rollouts of the official generator on the training inputs (the inference-time distribution), and (b) NCM-style uniform code corruption with level 0 preserved, where every code is replaced by a uniformly random code with probability $u\sim\mathcal{U}(0,1)$ drawn per sequence. We minimize
\begin{align}
\mathcal{L}_{\mathrm{APC}}(\omega)=\sum_{k}\lambda_k\Big(&\tfrac{1}{Ld}\big\|\zt^{(k)}-\zh^{(k)}_{\mathrm{GT}}\big\|_F^2\nonumber\\
&+\alpha\,\KL\big(D_\psi(\zh^{(k)}_{\mathrm{GT}})\,\big\|\,D_\psi(\zt^{(k)})\big)\Big), \label{eq:loss}
\end{align}
with $\lambda_k\propto 1/\mathbb{E}\|\zh^{(k)}_{\mathrm{GT}}\|^2$ normalized to sum to the number of levels, $\alpha=0.1$, and the KL averaged over positions. For decoders with large vocabularies (151{,}936 for the Qwen3 tokenizer used by WMT14, 50{,}257 for GPT-2 on XSum) the KL is evaluated on a random subset of (sequence, level) pairs per step and rescaled, which keeps the estimator unbiased. Rollouts are drawn with probability $0.75$ and corrupted trajectories otherwise. Target labels enter only~\eqref{eq:loss}; at inference $C_\omega$ receives $\zh_{\mathrm{pred}}$, the prefix and $k$.

\subsection{Rate-distortion Cauchy halting and token freezing}
For $k\ge k_{\min}$ define the normalized latent flux and the decoder stability
\begin{align}
\Delta_k&=\frac{\|\zt^{(k)}-\zt^{(k-1)}\|_F}{\|\zt^{(k-1)}\|_F+\epsilon},\label{eq:flux}\\
S_k&=\frac{1}{L}\sum_{t=1}^{L}\mathbb{I}\big[y^{(k)}_t=y^{(k-1)}_t\big],\quad y^{(k)}_t=\argmax D_\psi(\zt^{(k)})_t.\label{eq:stab}
\end{align}
A batch halts at the first $k^\ast\le K-2$ at which every sequence satisfies $\Delta_{k}<\tau$ (flux rule) or $1-S_k<\tau$ (stability rule); the output is $D_\psi(\zt^{(k^\ast)})$ and levels $k^\ast+1,\dots,K-1$ are never executed. Halting is applied per batch because only then is generator computation actually saved by the parallel implementation. For token freezing we compute a local flux $\delta_{k,t}=\|\zt^{(k)}_t-\zt^{(k-1)}_t\|_2/(\|\zt^{(k-1)}_t\|_2+\epsilon)$, area-pool it to the $\ell_{k+1}$ segments of the next level, and freeze segments with $\delta<\tau_{\mathrm{tok}}$: their latent rows are carried forward unchanged at every later level. Freezing does not reduce FLOPs in the dense block-diagonal implementation; we evaluate it for its effect on quality.

\subsection{Hyper-parameter selection without test leakage}
\label{sec:select}
All choices are made on the first $N_{\mathrm{tr}}$ items of each official test file, which are disjoint from the reported evaluation items. Within these, $80\%$ are used for gradient steps and $20\%$ as a validation set. (i) The training checkpoint with the lowest validation loss~\eqref{eq:loss} on rollout trajectories is kept, including the identity at initialization (epoch $-1$). (ii) The step-size scale $s$ is chosen from $\{0,0.25,0.5,1\}$ by \emph{closed-loop} validation quality (solve rate, BLEU or ROUGE-L); for the decode-path variant a separate $s_k$ is chosen per level using the level-$k$ decoding quality (cell accuracy for Sudoku). Ties are broken towards the smaller scale. (iii) The halting rule and $\tau$ are chosen on the same $N_{\mathrm{tr}}$ items as the setting with the fewest mean levels whose quality is within $0.5$ points of the full-depth model. Run 1 of our study (``in-loop, no selection'') omits (i) and (ii); it is reported as an ablation.





\section{Theoretical Analysis}
\label{sec:theory}
We analyze the drift recursion~\eqref{eq:drift} and the halting error. Norms are Frobenius norms on $\mathbb{R}^{L\times d}$, and expectations are over inputs and sampled codes. We index generation levels $k=1,\dots,K$ in this section for readability (level $0$ in code is $k=1$ here).

\begin{assumption}[Lipschitz upsampling]\label{as:lip}
$A_k\triangleq\Phi_k\circ\Up_k$ is linear with $\|A_k\|_{\mathrm{op}}\le L_p$.
\end{assumption}
\begin{assumption}[Drift coherence in expectation]\label{as:coh}
Let $\nu_{k-1}=e^{(k-1)}/\|e^{(k-1)}\|$ if $e^{(k-1)}\neq0$ and an arbitrary fixed unit vector otherwise. There is $c_0\in(0,1]$ such that $\mathbb{E}\big[\langle \nu_{k-1},\epsilon_k\rangle\big]\ge c_0\,\sigma_k$ for every $k$, where $\sigma_k\triangleq\mathbb{E}\|\epsilon_k\|>0$.
\end{assumption}
\begin{assumption}[$\epsilon_{\mathrm{corr}}$-accurate contractive corrector]\label{as:corr}
For every $k$ and every $z$ with $\|z-\zh^{(k)}_{\mathrm{GT}}\|\le R$, the corrected point $\mathcal{C}_k(z)=z-\eta_kC_\omega(z,h_{\mathrm{prefix}},k)$ satisfies $\|\mathcal{C}_k(z)-\zh^{(k)}_{\mathrm{GT}}\|\le(1-\gamma_k)\|z-\zh^{(k)}_{\mathrm{GT}}\|+\epsilon_{\mathrm{corr}}$ with $\gamma_k\in[0,1)$, and all iterates remain in the ball of radius $R$.
\end{assumption}

Assumption~\ref{as:lip} holds for the released tokenizers because linear interpolation and $\Phi_k$ are linear maps; $L_p\le(1-\varrho)+\varrho\|\mathrm{Conv1d}_k\|_{\mathrm{op}}$. Assumption~\ref{as:coh} states that, on average, the code error at level $k$ has a positive component along the drift accumulated so far, i.e., finer codebooks do not cancel it; it only constrains a first moment and is measured on the released checkpoints in Remark~\ref{rem:c0}. Assumption~\ref{as:corr} is the contraction property that the objective~\eqref{eq:loss} aims to induce.

\begin{proposition}[Lyapunov contraction under predictor-corrector revision]\label{prop:contraction}
Under Assumptions~\ref{as:lip}--\ref{as:coh}, purely additive MS-RVQ generation satisfies
\begin{equation}
\mathbb{E}\big\|\zh^{(K)}-\zh^{(K)}_{\mathrm{GT}}\big\|\ \ge\ c_0\sum_{k=1}^{K}\sigma_k=\Omega\Big(\sum_{k=1}^{K}\sigma_k\Big), \label{eq:lower}
\end{equation}
and, if $\#\{t:\hat r_{k,t}\neq r^\ast_{k,t}\}\le m_k$, then $\sigma_k\le L_p\,m_k\,\mathrm{diam}(\mathcal{V}_k)$. Under Assumption~\ref{as:corr}, with the in-loop update $\zt^{(k)}=\mathcal{C}_k\big(\zt^{(k-1)}+A_kE_k(\hat r_k)\big)$ and $\tilde e^{(k)}=\zt^{(k)}-\zh^{(k)}_{\mathrm{GT}}$, the Lyapunov function $V_k=\|\tilde e^{(k)}\|$ obeys $V_k\le(1-\gamma_k)(V_{k-1}+\|\epsilon_k\|)+\epsilon_{\mathrm{corr}}$, hence
\begin{equation}
\mathbb{E}V_K\le\sum_{k=1}^{K}\Big(\prod_{j=k}^{K}(1-\gamma_j)\Big)\sigma_k+\epsilon_{\mathrm{corr}}\sum_{k=1}^{K}\prod_{j=k+1}^{K}(1-\gamma_j). \label{eq:upper}
\end{equation}
If $\gamma_j\ge\gamma>0$ and $\sigma_k\le\bar\sigma$, then $\mathbb{E}V_K\le(\bar\sigma(1-\gamma)+\epsilon_{\mathrm{corr}})/\gamma$ uniformly in $K$, whereas~\eqref{eq:lower} grows linearly in $K$.
\end{proposition}
\begin{proof}[Proof sketch]
For the lower bound, Cauchy--Schwarz gives pointwise $\|e^{(k)}\|\ge\langle \nu_{k-1},e^{(k-1)}+\epsilon_k\rangle=\|e^{(k-1)}\|+\langle \nu_{k-1},\epsilon_k\rangle$. Taking expectations and applying Assumption~\ref{as:coh} gives $\mathbb{E}\|e^{(k)}\|\ge\mathbb{E}\|e^{(k-1)}\|+c_0\sigma_k$, which telescopes linearly; no Jensen step is needed. The bound on $\sigma_k$ follows from linearity of $A_k$: $\|\epsilon_k\|\le L_p\|E_k(\hat r_k)-E_k(r_k^\ast)\|\le L_p\,m_k\,\mathrm{diam}(\mathcal{V}_k)$. For the upper bound, $\tilde e^{(k)}=\mathcal{C}_k(\zt^{(k-1)}+A_kE_k(\hat r_k))-\zh^{(k)}_{\mathrm{GT}}$ and $\zt^{(k-1)}+A_kE_k(\hat r_k)-\zh^{(k)}_{\mathrm{GT}}=\tilde e^{(k-1)}+\epsilon_k$ by~\eqref{eq:cum}; Assumption~\ref{as:corr} and the triangle inequality give the one-step bound, which is unrolled by induction (Appendix~\ref{app:p1}). The uniform bound sums a geometric series.
\end{proof}

\begin{assumption}[Decoder regularity on the fine-drop support]\label{as:dec}
The logit map $f=D_\psi$ is $L_D$-Lipschitz, $\|f(a)-f(b)\|_F\le L_D\|a-b\|_F$, on the segment between $\zt^{(k^\ast)}$ and $\zt^{(K)}$, which lies in the support of the fine-drop training distribution~\eqref{eq:finedrop}.
\end{assumption}
\begin{assumption}[Geometric decay of increments]\label{as:geo}
The increments $\Delta\zt^{(j)}=\zt^{(j)}-\zt^{(j-1)}$ satisfy $\|\Delta\zt^{(j+1)}\|\le\rho\|\Delta\zt^{(j)}\|$ for $j\ge k^\ast$, with $\rho\in(0,1)$.
\end{assumption}

\paragraph{Off-manifold coupling.} Proposition~\ref{prop:contraction} treats $\sigma_k$ as given, but in closed loop the generator's error at level $k+1$ depends on its input. Let $\mathcal{M}_k=\{\sum_{j\le k}A_jE_j(r_j)\}$ be the lattice of level-$k$ codeword sums on which the generator was trained, and write $W_k=\|\zh^{(k)}-\zh^{(k)}_{\mathrm{GT}}\|$ for the error \emph{before} correction at level $k$.
\begin{assumption}[Off-manifold sensitivity]\label{as:offm}
There are $\sigma^{\mathrm{on}}_{k+1}$ and $L_G\ge0$ such that, for inputs $\zt$ of the generator at level $k+1$, $\mathbb{E}\|\epsilon_{k+1}\|\le\sigma^{\mathrm{on}}_{k+1}+L_G\,\mathbb{E}[\mathrm{dist}(\zt,\mathcal{M}_k)]$, with $\sigma^{\mathrm{on}}_{k+1}$ the error on lattice inputs.
\end{assumption}
\begin{corollary}[When in-loop correction helps]\label{cor:LG}
Suppose Assumptions~\ref{as:lip}, \ref{as:corr} and~\ref{as:offm} hold, and Assumption~\ref{as:coh} holds at level $k+1$ for the uncorrected process with error $\epsilon^{\mathrm{u}}_{k+1}$ and $\sigma^{\mathrm{u}}_{k+1}=\mathbb{E}\|\epsilon^{\mathrm{u}}_{k+1}\|$. Then one in-loop correction at level $k$ yields a smaller expected error after level $k+1$ than not correcting whenever
\begin{equation}
L_G<L_G^{\ast}\triangleq\frac{\gamma_k\,\mathbb{E}W_k-\epsilon_{\mathrm{corr}}-\big(\sigma^{\mathrm{on}}_{k+1}-c_0\,\sigma^{\mathrm{u}}_{k+1}\big)}{D_k},\qquad D_k=\mathbb{E}\big[\mathrm{dist}(\zt^{(k)},\mathcal{M}_k)\big]. \label{eq:LGthr}
\end{equation}
\end{corollary}
\begin{proof}
For the corrected process, the triangle inequality and Assumption~\ref{as:corr} give $\|\tilde e^{(k+1)}\|\le\|\tilde e^{(k)}\|+\|\epsilon_{k+1}\|\le(1-\gamma_k)W_k+\epsilon_{\mathrm{corr}}+\|\epsilon_{k+1}\|$; taking expectations and applying Assumption~\ref{as:offm}, $\mathbb{E}\|\tilde e^{(k+1)}\|\le(1-\gamma_k)\mathbb{E}W_k+\epsilon_{\mathrm{corr}}+\sigma^{\mathrm{on}}_{k+1}+L_GD_k$. For the uncorrected process, the pointwise argument of Proposition~\ref{prop:contraction} with Assumption~\ref{as:coh} gives $\mathbb{E}\|e^{(k+1)}\|\ge\mathbb{E}W_k+c_0\sigma^{\mathrm{u}}_{k+1}$. The first bound is below the second exactly when $L_GD_k<\gamma_k\mathbb{E}W_k-\epsilon_{\mathrm{corr}}-(\sigma^{\mathrm{on}}_{k+1}-c_0\sigma^{\mathrm{u}}_{k+1})$, which is~\eqref{eq:LGthr}.
\end{proof}
The threshold makes the trade-off explicit: the benefit of correction scales with the contraction $\gamma_k$ and the size of the accumulated error $W_k$, and the cost with the generator's off-manifold sensitivity $L_G$ times the lattice distance $D_k$ the correction creates. Because a contraction must move the latent by at least $\gamma_k W_k-\epsilon_{\mathrm{corr}}$, $D_k$ is typically of the same order as the benefit, so~\eqref{eq:LGthr} requires $L_G$ well below one unless the correction stays close to $\mathcal{M}_k$, which motivates projection (Section~\ref{sec:results}) and off-lattice robustness training (Section~\ref{sec:concl}).

\begin{proposition}[Rate-distortion bound for Cauchy halting]\label{prop:rd}
Suppose halting occurs at $k^\ast$ with $\|\Delta\zt^{(k^\ast)}\|\le\tau$ (absolute flux, i.e., $\Delta_{k^\ast}\|\zt^{(k^\ast-1)}\|\le\tau$). Under Assumptions~\ref{as:dec}--\ref{as:geo}, $\|\zt^{(K)}-\zt^{(k^\ast)}\|\le\sum_{j=k^\ast+1}^{K}\|\Delta\zt^{(j)}\|\le\rho\tau/(1-\rho)\le\tau/(1-\rho)$, and with $p_\psi(\cdot\mid z)=\mathrm{softmax}(f(z))$ applied per position,
\begin{equation}
\sum_{t=1}^{L}\KL\big(p_\psi(y_t\mid\zt^{(K)})\,\big\|\,p_\psi(y_t\mid\zt^{(k^\ast)})\big)\le\frac{L_D^2\tau^2}{2(1-\rho)^2}. \label{eq:rd}
\end{equation}
\end{proposition}
\begin{proof}[Proof sketch]
The first chain is the triangle inequality and a geometric series. For two logit vectors $a,b\in\mathbb{R}^{V}$, $\KL(\mathrm{softmax}(a)\|\mathrm{softmax}(b))$ is the Bregman divergence of the log-partition function $\mathrm{lse}$, whose Hessian $\mathrm{diag}(p)-pp^\top$ has quadratic form $\mathrm{Var}_p(h)\le\|h\|_\infty^2\le\|h\|_2^2$; Taylor's theorem with integral remainder yields $\KL\le\frac12\|a-b\|_2^2$. Summing over positions and applying Assumption~\ref{as:dec} gives~\eqref{eq:rd} (Appendix~\ref{app:p2}).
\end{proof}

\begin{remark}[Role of fine-drop]
Proposition~\ref{prop:rd} compares two decoder outputs; it does not by itself say that $D_\psi(\zt^{(k^\ast)})$ is a sensible output. Fine-drop training places $\zt^{(k^\ast)}$ on the decoder's training support, so the Lipschitz constant in Assumption~\ref{as:dec} is a local property of a region the decoder was fitted on rather than an extrapolation. If $D_\psi(\zt^{(K)})$ is correct with margin larger than $\sqrt{2\cdot\mathrm{RHS}}$ in total variation at every position (Pinsker's inequality~\citep{pinsker1964,cover2006elements}), the halted argmax output equals the full-depth one.
\end{remark}


\begin{remark}[Measured drift coherence and off-manifold sensitivity]\label{rem:c0}
On greedy rollouts of the released checkpoints (64 Sudoku-hard boards, 12 WMT14 sentences), the per-level ratio $\hat c_{0,k}=\mathbb{E}\langle \nu_{k-1},\epsilon_k\rangle/\sigma_k$ has mean \CzeroSud{} on Sudoku-hard (negative at 8 of 11 levels) and \CzeroWmt{} on WMT14 (negative at 1 of 10). Assumption~\ref{as:coh} therefore fails on Sudoku, where finer levels on average partially \emph{cancel} accumulated drift, and holds only weakly ($c_0\approx0.08$) on WMT14. The linear lower bound~\eqref{eq:lower} is thus a worst-case statement that the released models do not attain, and additive drift is not the bottleneck it was hypothesized to be; this leaves a corrector little error to remove. Perturbing rollout states off $\mathcal{M}_k$ with the in-loop corrector at scales $s\in\{0,0.25,0.5,1\}$ and regressing next-level code error on lattice distance within each level gives a median slope $\hat L_G=$ \LGSud{} (Sudoku-hard) and \LGWmt{} (WMT14). On Sudoku-hard the next-level error rises from \SigOnSud{} at $s=0$ to \SigOneSud{} at $s=1$ as the mean lattice distance grows from \DistZeroSud{} to \DistOneSud; on WMT14 the one-step error is almost unchanged (\SigOnWmt{} to \SigOneWmt), so the closed-loop collapse there comes from compounding of off-lattice states across levels rather than from a single step.
\end{remark}

\begin{remark}[Empirical status of the assumptions]
We estimated the constants on reference trajectories of the released tokenizers (Section~\ref{sec:diag}). The empirical decoder ratio $\hat L_D^2=\max\KL/(\frac12\|\zh^{(K)}-\zh^{(k)}\|^2)$ is small ($\le\LDSud$ on Sudoku-hard and $\le\LDWmt$ on WMT14), so Assumption~\ref{as:dec} is benign. Assumption~\ref{as:geo}, however, is violated: the median ratio $\|\Delta\zh^{(j+1)}\|/\|\Delta\zh^{(j)}\|$ exceeds 1 at \RhoGtOneSud{} of 10 Sudoku-hard level pairs and \RhoGtOneWmt{} of 9 WMT14 pairs, and the relative flux at the second-to-last level is \FluxLastSud{} (Sudoku-hard) and \FluxLastWmt{} (WMT14), far above any useful $\tau$. Proposition~\ref{prop:rd} is therefore vacuous for these models, which is consistent with the empirical halting results.
\end{remark}
